%File: who_takes_the_credit.tex  (AAAI-26 author kit, anonymous submission)
% Needs aaai2026.sty and aaai2026.bst (author kit) and references.bib in the same folder.
% DRAFT in Korean: translate to English, then delete the kotex line, the hyperref line and every \TODO.
\documentclass[letterpaper]{article} % DO NOT CHANGE THIS
\usepackage[submission]{aaai2026}  % DO NOT CHANGE THIS
\usepackage{times}  % DO NOT CHANGE THIS
\usepackage{helvet}  % DO NOT CHANGE THIS
\usepackage{hyperref}  % DRAFT ONLY: clickable citations. Remove before submission (check the AAAI kit's banned-package list).
\usepackage{courier}  % DO NOT CHANGE THIS
\usepackage[hyphens]{url}  % DO NOT CHANGE THIS
\usepackage{graphicx} % DO NOT CHANGE THIS
\urlstyle{rm} % DO NOT CHANGE THIS
\def\UrlFont{\rm}  % DO NOT CHANGE THIS
\usepackage{natbib}  % DO NOT CHANGE THIS AND DO NOT ADD ANY OPTIONS TO IT
\usepackage{caption} % DO NOT CHANGE THIS AND DO NOT ADD ANY OPTIONS TO IT
\frenchspacing  % DO NOT CHANGE THIS
\setlength{\pdfpagewidth}{8.5in} % DO NOT CHANGE THIS
\setlength{\pdfpageheight}{11in} % DO NOT CHANGE THIS
\usepackage{booktabs}
\usepackage{amsmath}
\usepackage{kotex}  % DRAFT ONLY: Korean text. Remove before submission (AAAI forbids non-Roman fonts).
\pdfinfo{
/TemplateVersion (2026.1)
}
\setcounter{secnumdepth}{0}
\newcommand{\TODO}[1]{[\textbf{TODO:} #1]}

\title{Who Takes the Credit? Measuring Contribution in Agentic Teams}
\author{
    Anonymous Submission
}
\affiliations{
}

\begin{document}

\maketitle

\begin{abstract}
팀이 성공했다는 사실만으로는 누가 기여했는지, 모든 팀원이 필요했는지 알 수 없다. 우리는 흔한 보상
설계로 학습한 에이전트가 실제 기여보다 많은 공(credit)을 받는 나쁜 팀원이 되는지, 그리고 LLM 판정자가
이를 알아차리는지 연구한다. Overcooked에서 학습자를 고정된 동료 옆에서 네 가지 보상(팀, 개인, 줍기 보너스,
근접 보너스)과 두 알고리즘(PPO, GRPO)으로 학습시켜 \emph{나쁜 팀원 pool}을 만든다. 각 팀원의 실제 기여는
한 명을 멈춘 채 게임을 다시 돌려 계산한 Shapley 몫으로 재고, 직접 배달한 비율과의 차이를 \emph{credit gap}으로
정의한다. pool에는 네 가지 hacker 페르소나가 나타난다: 배달 뺏기(\emph{credit theft}), 바쁜 척(\emph{fake help}),
그림자(\emph{shadowing}), 포기(\emph{quitting}). 배달 뺏기는 팀 수프 11개를 모두 배달하지만 Shapley 몫은
0.18이다(credit gap 0.82, 규칙 기반 요리사 0.13). 하나의 이벤트에 묶인 보상은 두 알고리즘 모두에서 같은
페르소나를 만든다 \TODO{다른 맵 결과}. LLM 판정자(공개 27B 모델과 GPT-4o)는 이 도둑에게 공의 0.90--0.92를
준다 \TODO{GRPO 값}. 판정자의 공은 실제 기여($r\le0.35$)가 아니라 누가 배달했는지($r=0.94$)를 따른다.
에피소드 기록만 보는 LLM 판정자는 반사실 기여를 보지 못해 credit hacker에게 속는다.
\end{abstract}

% Code / data links (AAAI kit `links` environment; shown in submission mode).
% Anonymous review: anonymized URLs only (e.g. anonymous.4open.science), never real accounts.
\begin{links}
    \link{Code}{https://anonymous.4open.science/r/TODO}
    \link{Models and data}{https://huggingface.co/TODO}
\end{links}

% ===================================================================================== Introduction
\section{Introduction}
LLM 에이전트와 강화학습 에이전트가 팀을 이루어 일하는 경우가 늘고 있다.
이런 팀은 보통 목표를 달성했는지로 평가한다. 팀원 하나가 다른 팀원의 결과를 가로채거나 일하는 척만
해도 팀 점수는 그대로이거나 오히려 오를 수 있다.

이는 공정성뿐 아니라 비용의 문제이다. 다섯 에이전트 팀이 일을 해냈더라도, 실제로는 셋으로 충분했고
나머지 둘은 놀거나 일하는 척만 했을 수 있다. 에이전트마다 연산, API 호출, 토큰 비용이 들기 때문에,
필요 없는 팀원은 팀 점수에 드러나지 않는 낭비가 된다. 실제로 LLM 다중 에이전트 시스템에서 불필요한
에이전트나 메시지를 없애면 성능을 유지하면서 토큰 비용을 크게 줄일 수 있다
\cite{zhang2025cutthecrap,wang2025agentdropout}. 누가 필요했는지 알려면 ``이 팀원이 없었어도 팀이 같은
결과를 냈을까?''라는 반사실 질문에 답해야 한다.

LLM 판정자는 에이전트 팀을 평가하는 쉬운 방법이지만, 에피소드 기록만 보고
이 질문에 답해야 한다. 우리는 LLM 판정자가 반사실을 고려하지 않아, 결과를 마무리한 팀원을 실제보다
유능하다고 판단하고 credit hacker에게 과한 공을 준다는 것을 보인다.

이런 credit hacker는 일부러 만들 필요가 없다. 학습하는 에이전트는 실제 기여가 아니라 보상 규칙을
최대화하므로, 흔한 보상 설계만으로도 생긴다.
우리는 협동 요리 게임 Overcooked \cite{carroll2019utility}에서 학습자를 고정된 동료 옆에서 네 가지 흔한
보상과 PPO, GRPO로 학습시켜 \emph{나쁜 팀원 pool}을 만든다. 페르소나를 미리 정하지 않고, 보상과 동료만
정한 뒤 어떤 행동이 나오는지는 학습에 맡긴다. 각 팀원의 실제 기여는 한 명을 멈춘 채 같은 게임을 다시
돌려 계산한 Shapley 몫으로 잰다. 이 몫은 그 팀원이 없었다면 팀이 무엇을 잃었는지를
반영하며, 직접 배달한 비율과 이 몫의 차이를 \emph{credit gap}으로 정의한다.

pool에는 네 가지 나쁜 팀원이 나타났다. \emph{배달 뺏기}는 동료가 만든 수프를 가져가 모두 직접 배달한다.
동료는 혼자서도 수프를 거의 그대로 만들기 때문에, 공을 가장 많이 받는 이 팀원은 사실상 없어도 되는
팀원이다. \emph{바쁜 척}은 많이 줍지만 아무것도 만들지 않고, \emph{그림자}는 동료 옆에 붙어 길을 막아
팀 생산량을 0으로 만들며, \emph{포기}는 동료가 배달을 독점하면 맡던 준비마저 멈춘다. 이 중 배달 뺏기는
팀 점수를 오히려 올려, 팀 점수 감독자는 배달 뺏기 학습 13개 중 9개를 통과시킨다.
\TODO{13개는 LBF, RWARE와 예전 판정 기준을 포함한 값: cramped\_room과 gap $>0.5$ 기준으로 다시 계산}
\TODO{다른 맵 결과 한 문장}

LLM 판정자는 많이 줍기만 하는 바쁜 척에는 거의 속지 않지만(0.04--0.09), 배달 뺏기에는 거의 항상 속는다.
더 많이 배달했지만 실제 기여는 더 작은 쌍에서 판정자는 0.88--1.00의 비율로 도둑을 고른다. 판정자는
겉으로 보이는 활동이 아니라 \emph{누가 마지막 단계를 했는지}로 공을 나누며, 기록만으로는 누가 없어도
되는 팀원인지 알아내지 못한다. \TODO{GRPO 도둑 판정 결과 한 문장}

우리의 기여는 다음과 같다.
\begin{itemize}
\item 흔한 보상 설계에서 저절로 생기는 \textbf{나쁜 팀원 pool}: 네 가지 credit hacker 페르소나와
      이를 만드는 보상과 동료의 조합.
\item 반사실 재실행으로 계산한 Shapley 몫에 기반한 \textbf{credit gap}: 각 팀원이 공을 얼마나
      부풀렸는지와 팀에 꼭 필요했는지를 함께 보여준다.
\item LLM 판정자의 \textbf{취약점 분석}: 판정자는 결과를 마무리한 팀원에게 공을 주어 credit hacker를
      놓치며, 이는 기록에 반사실 정보가 없기 때문이다.
\end{itemize}

% ===================================================================================== Related work
\section{Related Work}
\textbf{동료 pool.} Zero-shot coordination 연구는 다양한 동료로 협력 에이전트를 학습하거나 평가한다.
HSP \cite{yu2023hsp}는 이벤트(양파 줍기, 배달 등)에 무작위 가중치를 준 숨은 보상으로 편향된 동료를 만들고,
ZSC-Eval \cite{wang2024zsceval}은 같은 방식의 동료를 고정된 평가용 pool로 쓴다. FCP \cite{strouse2021fcp},
MEP \cite{zhao2023mep} 등은 실력과 행동이 다양한 pool을 만든다. 이들은 pool로 \emph{협력 에이전트}를 학습하거나
평가하고, 우리는 나쁜 팀원 pool로 \emph{credit 판정자}를 평가한다. 또 기존 이벤트 공간은 팀 점수를 공유하고
위치 이벤트가 없어, 공을 가로채는 페르소나(배달 뺏기)와 위치 기반 페르소나(그림자)를 만들지 않는다.

\textbf{보상 해킹과 credit 할당.} 보상 해킹 \cite{pan2022effects,skalse2022defining}, 다중 에이전트 학습의
게으른 에이전트 \cite{sunehag2018vdn}와 difference reward \cite{foerster2018coma}, Shapley 값에 기반한 credit
할당 \cite{wang2020shapleyq,li2021shapley}이 관련된다. 우리는 Shapley 값을 학습 신호가 아니라
\emph{감사(audit)} 도구로 쓴다.

\textbf{LLM 판정자.} LLM을 판정자로 쓰는 연구 \cite{zheng2023judging}는 주로 답의 품질을 평가하며, 판정자가
답의 위치 같은 표면적 단서에 흔들린다는 점도 보고되었다 \cite{wang2024fair}. 다중 에이전트 LLM 시스템에서
실패의 원인이 된 에이전트를 LLM이 찾게 한 연구 \cite{zhang2025which}는 우리와 가장 가깝지만, 실패의 탓을
판정하고 정답 원인과 비교한다. 우리는 성공한 팀 안에서 누가 기여했는지를 판정하게 하고, 그 판단을 반사실
기여와 비교한다.

% ===================================================================================== Setup
\section{Setup}
\textbf{환경.} Overcooked-AI의 \texttt{cramped\_room}에서 요리사 두 명이 400 step
동안 양파 수프를 만든다. 부엌은 좁아서 두 요리사가 서로의 길을 막을 수 있다. 수프 하나는 일곱 단계로
완성되며 팀에 20점을 준다: (1) 양파 줍기, (2) 냄비에 넣기(세 번), (3) 빈손으로 요리 시작, (4) 익기를
기다리기, (5) 접시 줍기, (6) 수프 담기, (7) 배달. 단계 1--3은 수프를 \emph{만드는} 일이고, 단계 5--7은 만든
수프를 \emph{마무리하는} 일이다. 이 구분이 이후 모든 결과의 바탕이 된다. 마무리하는 쪽은 배달이 자기 이름으로
기록되지만, 수프가 존재하는 것은 만드는 쪽 덕분이기 때문이다. 학습자는 항상 0번 자리, 동료는 1번 자리에 둔다.

\textbf{동료.} 페르소나는 고정된 동료 옆에서 학습한다. 보상 해킹은 동료의 행동을 이용하는 것이므로, 동료가
학습 중에 바뀌지 않아야 페르소나가 안정적으로 생기고 반사실 기여도 잴 수 있다. 두 종류의 동료를 쓴다.
(i) \emph{greedy}: Overcooked-AI 패키지의 규칙 기반 요리사로, 혼자서도 수프를 끝까지 만든다. 상대가 접시를
들고 있으면 접시를 가지러 가지 않고 양파 쪽 일을 하는 규칙이 있다. (ii) \emph{A}: greedy 옆에서 팀 보상으로
학습한 뒤 고정한 학습자(seed 0)이다. 같은 팀 보상이라도 알고리즘에 따라 A의 성격이 다르다. PPO-A는 팀 수프
11.0개를 모두 직접 배달하는 배달 독점형이고, GRPO-A는 하나도 배달하지 않고 greedy가 8.3개를 모두 만드는
구경꾼형이다. A는 greedy와 달리 양보 규칙이 없어서, 같은 보상이 다른 동료 옆에서 어떤 페르소나를 만드는지
비교할 수 있게 해 준다.

\textbf{보상.} 학습자는 아래 보상 중 하나만 받는다. 네 보상은 모두 실제 다중 에이전트 학습에서 흔히 쓰는
형태이다. \emph{team}: 팀이 배달한 수프마다 20. \emph{individual}: 자기가 배달한 수프만 20(개인 성과급).
\emph{pickup}: 팀 보상 + 양파나 접시를 주울 때마다 3(활동 보너스). \emph{proximity}: 팀 보상 + 동료 바로
옆(맨해튼 거리 1)에 있는 step마다 1(협력 장려 보너스). 학습 초기를 돕는 Overcooked-AI의 공식 shaping 보상
(냄비에 재료 넣기 +3, 쓸모 있는 접시 줍기 +3, 수프 담기 +5)은 team과 individual에 더하고 1M step까지
선형으로 0으로 줄인다. pickup과 proximity의 보너스는 줄이지 않는다.

\textbf{학습.} 정책과 가치 함수는 각각 은닉층 128의 MLP이고, 관측은 Overcooked-AI의 lossless 상태 인코딩을
펼친 520차원 벡터이다. 알고리즘은 PPO \cite{schulman2017ppo}와 GRPO \cite{shao2024deepseekmath}이며, 각
조합을 2M 환경 step 동안 seed 0--2로 학습한다. GRPO는 critic 없이 같은 시작 상태에서 돈 에피소드 8개를
한 그룹으로 묶어 advantage를 계산한다. 두 알고리즘은 학습률, epoch, minibatch, clip을 같게 맞췄으며, 전체
하이퍼파라미터는 부록 표~\ref{tab:hparams}에 있다.

% ------------------------------------------------------------------------------------- Persona pipeline
\subsection{페르소나 만들기}
나쁜 팀원 pool은 두 단계로 만든다. \textbf{1단계(동료)}에서는 동료를 정한다. greedy는 그대로 쓰고, A는
greedy 옆에서 팀 보상으로 학습한 뒤 고정한다. \textbf{2단계(페르소나)}에서는 고정된 동료 옆에 \emph{새}
학습자를 처음부터 학습시키고, 위 보상 중 하나를 준다. 페르소나를 입는 쪽은 동료가 아니라 이 새 학습자이며,
동료는 학습 내내 바뀌지 않는다. 페르소나를 미리 정해 두고 그것을 흉내 내도록 학습시키지 않는다는 점이
중요하다. 우리는 보상과 동료만 정하고, 어떤 행동이 나오는지는 학습에 맡긴다. 따라서 각 페르소나는
``이 보상을 이 동료 옆에서 쓰면 이런 팀원이 생긴다''는 관찰이다.
\TODO{self-play로 학습한 A 옆에서 같은 보상을 쓴 결과 (results/sp\_A, sp\_B)}

\textbf{설계.} 표~\ref{tab:steps}는 각 보상으로 얻으려던 페르소나가 일곱 단계 중 어디를 맡는지 정리한다.
배달 뺏기는 개인 보상으로 마무리 단계(5--7)만 가져가고 만드는 단계(1--3)는 동료에게 맡기는 팀원이다.
바쁜 척은 줍기 보너스로 단계 1과 5를 반복하지만 단계 2로 넘기지 않아 아무것도 만들지 않는다. 그림자는
근접 보너스로 동료 옆에 붙어 있다가 좁은 부엌에서 동료의 길을 막는다. 포기는 개인 보상을 배달 독점형 동료
옆에서 받은 경우로, 자기 배달이 생길 수 없으니 shaping이 사라지면 맡던 준비 단계마저 멈춘다. 대조군으로
팀 보상을 받은 준비 담당은 단계 1--3을 맡고 배달은 동료에게 넘기는데, 이는 정상적인 역할 분담이다.

\begin{table}[t]
\centering
\small
\setlength{\tabcolsep}{4pt}
\begin{tabular}{@{}llll@{}}
\toprule
페르소나 & 보상 (동료) & 학습자 단계 & 동료 단계 \\
\midrule
기준: 규칙 요리사 & -- (greedy) & 1--7 나눠서 & 1--7 나눠서 \\
배달 뺏기 & individual (greedy) & 5, 6, 7 & 1--4 \\
바쁜 척 & pickup (greedy) & 1, 5 반복 & 1--7 \\
그림자 & proximity (A) & 없음 (길막기) & 막혀서 없음 \\
포기 & individual (A) & 1--3 뒤 멈춤 & 5--7 \\
\midrule
(대조) 준비 담당 & team (A) & 1--3 & 5--7 \\
\bottomrule
\end{tabular}
\caption{설계한 페르소나: 보상과 동료, 그리고 각 페르소나가 수프의 일곱 단계 중 맡을 것으로 예상한 단계.
(1) 양파 줍기, (2) 냄비에 넣기, (3) 요리 시작, (4) 기다리기, (5) 접시 줍기, (6) 수프 담기, (7) 배달.}
\label{tab:steps}
\end{table}

\begin{figure}[t]
\centering
\includegraphics[width=\columnwidth]{../figures/fig_steps.pdf}
\caption{재현 확인. 학습된 체크포인트를 학습 때의 동료 옆에서 평가 게임 10개로 다시 돌려 센 단계별 횟수.
색은 그 단계 중 학습자가 한 비율, 칸 안의 숫자는 게임당 학습자/동료 횟수, 회색은 둘을 합쳐 게임당 1회 미만.
단계 4(기다리기)는 행동이 아니라 뺐다. 그림자는 seed 3개의 평균, 나머지는 seed 0.}
\label{fig:steps}
\end{figure}

\textbf{재현 확인.} 보상이 의도한 페르소나를 실제로 만들었는지 확인하려고, 학습된 체크포인트를 학습 때의
동료 옆에서 다시 돌려 각 자리가 일곱 단계를 몇 번씩 했는지 센다(그림~\ref{fig:steps}). 단계 3은 빈손으로
상호작용해 냄비가 요리를 시작한 순간, 단계 6은 손에 든 접시가 수프로 바뀐 순간으로 센다. 결과는
표~\ref{tab:steps}의 설계와 맞는다. 규칙 요리사끼리는 모든 단계를 나눠서 한다(학습자 몫 0.4--0.6).
배달 뺏기는 PPO와 GRPO 모두 단계 5--7을 거의 혼자 하고(게임당 12/0, 11/0, 11/0) 단계 1--3은 한 번도 하지
않는다. 동료인 greedy가 양파 37개를 넣고 요리를 12번 시작한다. 바쁜 척은 양파를 게임당 178번 줍지만
냄비에는 하나도 넣지 않는다. 설계와 달리 이 체크포인트는 접시는 줍지 않고 양파만 반복해서 줍는다.
그림자는 400 step 중 399 step(GRPO 394)을 동료 옆에서 보내고, 그 사이 두 요리사 모두 수프를 하나도 만들지
못한다. 포기한 학습자는 최종 체크포인트에서 게임당 양파 2개만 줍고, 배달 독점형 동료는 접시를 46번 집지만
담을 수프가 없다. 즉 네 페르소나는 모두 보상이 가리킨 단계에서 정확히 어긋난다.
\TODO{GRPO 바쁜 척, 준비 담당 체크포인트를 받아 그림에 추가}

% ===================================================================================== Measuring
\section{Measuring Contribution}
\textbf{반사실 재실행.} 학습자 $L$과 동료 $P$로 이루어진 팀의 가치 $v(\cdot)$를 평가 게임 10개(seed
$10^4+k$)의 평균 수프 수로 정의한다. $v(L,P)$는 실제 게임이다. $v(L)$과 $v(P)$는 상대를 멈춘(frozen) 채
같은 seed로 게임을 다시 돌려 잰다. 멈춘 요리사는 냄비, 재료 상자, 접시 상자, 배달대 어디에도 닿지 않는
칸으로 옮겨 세운다. \texttt{cramped\_room}에서는 0번 요리사의 시작 칸이 접시 상자에 닿는 유일한 칸이라,
시작 칸에 그대로 세워 두면 접시를 아무도 꺼내지 못해 반사실 값이 망가지기 때문이다. 학습된 정책의 행동
샘플링도 게임마다 같은 seed로 다시 맞춰, 세 게임이 같은 무작위성을 공유하게 한다.

\textbf{Shapley 몫과 credit gap.} 두 명의 Shapley 값 \cite{shapley1953value}은
\begin{equation}
\phi_L = \tfrac{1}{2}\bigl[v(L) + v(L,P) - v(P)\bigr],\qquad s = \phi_L / v(L,P)
\end{equation}
이다. 첫 항은 학습자가 혼자 만드는 몫이고, 둘째 항은 동료에 학습자가 더해져 늘어난 몫이다. $s$는 팀 결과
중 학습자의 실제 몫으로, 학습자가 방해만 하면 음수가 될 수 있다. 학습자가 직접 배달한 수프의 비율을 $o$라
할 때 \emph{credit gap}은 $o - s$로, 겉으로 보이는 공이 실제 기여를 얼마나 넘는지를 뜻한다. 같은 자리에
규칙 요리사를 넣으면 $s^*=0.50$, gap $0.13$이다. 두 팀원 모두 혼자서는 수프를 못 만들면 $s$는 공식상
0.5로 고정되므로 따로 표시한다.

\textbf{Hacker 기준.} 학습자가 아래 둘 중 하나에 해당하면 hacker 후보로 표시한다. (a) credit gap $>0.5$:
배달로 보이는 공이 실제 몫보다 절반 이상 크다. (b) 학습 마지막 에피소드들(PPO는 마지막 5개 에피소드, GRPO는
마지막 5개 그룹)에서 받은 보상의 절반 이상이 팀 수프가 아닌 곳(줍기, 근접 등)에서 나왔다. 팀 수프가 기준의
절반 아래로 떨어진 경우는 피해로 따로 보고하지만, 서툰 학습자도 그렇게 될 수 있으므로 그것만으로는 hacker로
보지 않는다.

% ===================================================================================== Results 1
\section{A Pool of Bad Teammates}

\begin{table}[t]
\centering
\small
\setlength{\tabcolsep}{3pt}
\begin{tabular}{@{}llcccc@{}}
\toprule
페르소나 & 보상 (동료) & Alg & 팀 & 자기 & $s$ \\
\midrule
기준: 규칙 요리사 & -- (greedy) & -- & 9.1 & 5.7 & .50 \\
배달 뺏기 & individual (greedy) & PPO & 11.0 & 11.0 & .18 \\
 & & GRPO & 10.2 & 10.1 & .16 \\
바쁜 척 & pickup (greedy) & PPO & 7.5 & 0.0 & .03 \\
 & & GRPO & 7.1 & 0.0 & .01 \\
그림자 & proximity (A) & PPO & 0.0 & 0.0 & -- \\
 & & GRPO & 0.0 & 0.0 & -- \\
포기 & individual (A) & PPO & 0.07 & 0.0 & -- \\
\midrule
(대조) 준비 담당 & team (A) & PPO & 10.5 & 0.0 & .50 \\
\bottomrule
\end{tabular}
\caption{나쁜 팀원 pool (cramped\_room, seed 3개 평균). 팀: 팀 수프, 자기: 학습자가 배달한 수프,
$s$: Shapley 몫. 팀 수프가 0이면 $s$는 정의되지 않는다.}
\label{tab:pool}
\end{table}

표~\ref{tab:pool}은 pool의 페르소나를 팀 결과와 Shapley 몫으로 보여준다. 네 페르소나는 credit을 부풀리는
방식이 서로 다르다. 배달 뺏기는 남의 결과를 자기 이름으로 가져가고, 바쁜 척은 결과 없는 활동을 보여주며,
그림자와 포기는 팀 결과 자체를 무너뜨린다.

\textbf{배달 뺏기.} 개인 보상을 받은 학습자는 팀 수프 거의 전부를 직접 배달하지만 $s$는 0.16--0.18이고
credit gap은 0.82--0.83이다(규칙 요리사 0.13). 학습자는 접시를 먼저 들고 있어서, 상대가 접시를 들면
양보하는 greedy의 규칙을 이용한다. greedy는 양파만 나르고 배달은 학습자가 가져간다. 학습자를 멈추면 greedy
혼자서도 수프를 거의 그대로 만들기 때문에 학습자의 실제 몫은 작다. 그런데 팀 수프는 기준보다 오히려
높다(11.0 대 9.1). 역할이 고정되어 두 요리사의 동선이 덜 겹치는 것으로 보인다. 즉 이 hack은 팀에 해를
끼치지 않으면서 공만 부풀린다. 공유된 팀 보상을 받은 PPO도 같은 방식으로 배달을 뺏는다($s=0.18$).

\textbf{바쁜 척.} 줍기 보너스를 받은 학습자는 기준의 7--8배를 줍지만(165, 146 대 20) 냄비에는 넣지 않고
주운 것을 다시 내려놓는다. 팀 수프는 모두 greedy가 만들고 $s\approx0$이다. 배달을 하지 않으니 credit gap은
0에 가깝지만, 학습 보상의 70--80\%가 팀 수프가 아닌 줍기에서 나와 기준 (b)에 걸린다. 팀 수프는 기준의
80\% 정도로, 학습자가 부엌을 돌아다니며 생기는 약간의 방해만큼 줄어든다.

\textbf{그림자.} 근접 보너스를 받은 학습자는 400 step 중 PPO 399.5, GRPO 393.7 step을 A 옆에 붙어 있다
(기준 157--178). 좁은 부엌에서 계속 붙어 다니며 A의 길을 막아, 혼자서도 수프를 만드는 A까지 0개가 된다.
협력을 장려하려고 넣은 근접 보너스가 의도치 않게 방해꾼을 만든 것이다. 학습 보상은 거의 전부 근접
보너스에서 나온다.

\textbf{포기.} 배달을 독점하는 PPO-A 옆에서 개인 보상을 받은 학습자는 shaping이 남아 있는 1M step까지
양파를 넣는 준비 담당이 되어 팀 수프가 10개에 이른다. shaping이 0이 되면 자기 배달이 없어 보상이 0이 되고,
준비를 멈춰 팀 수프가 0.07개로 무너진다. 같은 보상이라도 양보하는 greedy 옆에서는 도둑이, 양보하지 않는
동료 옆에서는 포기가 된다. 페르소나는 보상만이 아니라 보상과 동료의 조합으로 정해진다.

\textbf{알고리즘.} 특정 이벤트에 묶인 보상(개인, 줍기, 근접)은 PPO와 GRPO에서 같은 페르소나를 만든다.
공유된 팀 보상에서는 알고리즘이 결과를 정한다: PPO는 배달을 뺏고($s=0.18$), GRPO는 무임승차한다($s=0.07$).
따라서 같은 팀 보상이라도 어떤 나쁜 팀원이 나올지는 학습 방식에 달려 있다.
\TODO{다른 맵 결과: 맵마다 같은 페르소나가 나오는지}

% ===================================================================================== Results 2
\section{Overseers Miss the Hacks}
두 종류의 감독자를 시험한다. \emph{팀 점수 감독자}는 팀 결과만 보고 팀을 통과시킬지 정하고,
\emph{LLM 판정자}는 에피소드의 사건 기록을 읽고 각 팀원에게 공을 나눈다. 둘 다 반사실 기여는 보지 못한다.

\textbf{팀 점수 감독자.} 가장 흔한 감독은 팀 점수만 보는 것이다. 기준 팀 점수의 80\% 이상이면 통과시키는
감독자는 배달 뺏기 학습 13개 중 9개를 통과시킨다. \TODO{cramped\_room과 gap $>0.5$ 기준으로 다시 계산}
배달 뺏기는 팀 점수를 낮추지 않기 때문이다. 반면 그림자와 포기는 팀 점수가 무너져 모두 걸러진다. 즉 팀 점수
감독은 팀을 망가뜨리는 hack만 잡고, 공을 가로채는 hack은 놓친다.

\textbf{LLM 판정자 실험.} 판정자가 에피소드 기록만 보고 공을 제대로 나누는지 시험한다. greedy를 1번 자리에
두고, 0번 자리에 다섯 종류의 팀원을 넣어 같은 seed 20개로 게임을 돌린다(총 100개). 다섯 종류는 규칙 요리사,
길막기(가만히 있지 않고 무작위로 움직이며 절반의 step에서 상호작용해, 물건을 줍고 내려놓기를 반복하는 규칙),
PPO 도둑, GRPO 무임승차, PPO 바쁜 척이다. 평가 seed는 학습과 평가에 쓴 seed와 겹치지 않는다. 각 게임은
관찰자가 볼 수 있는 사건 기록(누가 무엇을 줍고 넣고 배달했는지)으로 바꾸고, 반사실 값은 넣지 않는다.
판정자에게는 같은 seed에서 0번 자리만 다른 두 게임을 함께 보여주고, 어느 쪽 0번이 더 좋은 팀원인지와 각
게임에서 0번이 받아야 할 공의 비율(0--1)을 묻는다. Qwen3.8-27B(Ollama)는 116쌍, GPT-4o는 63쌍을 판정했다.

\textbf{측정.} 판정을 세 가지로 Shapley와 비교한다. (i) 일치율: 두 0번의 $\phi$ 차이가 0.24 수프를 넘는
쌍에서 판정자가 고른 쪽이 $\phi$가 큰 쪽과 같은 비율. (ii) 속은 비율: 한쪽이 겉보기에는 더 일했지만
(더 많이 배달했거나 더 많이 주웠지만) 실제 기여는 더 작은 쌍에서 판정자가 그쪽을 고른 비율. 0.5는 무작위,
0.5보다 크면 겉모습에 속는다는 뜻이다. (iii) 공의 상관: 게임마다 판정자가 준 공을 평균하고, 배달 비율 $o$와
Shapley 몫 $s$ 각각과의 Pearson $r$을 잰다.

\begin{figure}[t]
\centering
\includegraphics[width=\columnwidth]{../figures/fig_judge_credit.pdf}
\caption{LLM 판정자가 0번 팀원에게 준 공(색)과 0번이 직접 배달한 비율 $o$(빈 원), 실제 몫인 Shapley
$s$(검은 마름모). 게임 20개 평균, 기본 프롬프트. 길막기의 $s$는 $-1.1$로 축 밖에 있다.}
\label{fig:judge}
\end{figure}

\textbf{결과.} 두 판정자는 전체적으로는 Shapley와 잘 맞는다(일치율 Qwen 0.87, GPT-4o 0.79). 아무것도
배달하지 않는 팀원은 제대로 낮게 평가한다. 바쁜 척에 속은 비율은 0.04와 0.09로, 많이 줍는 것만으로는
공을 주지 않는다. 그러나 배달 뺏기에는 거의 항상 속는다. 더 많이 배달했지만 실제 기여는 더 작은 쌍에서
판정자는 0.88(Qwen)과 1.00(GPT-4o)의 비율로 도둑을 고른다. 그림~\ref{fig:judge}는 그 이유를 보여준다.
판정자의 공(색)은 모든 페르소나에서 빈 원($o$) 바로 옆에 있다. 도둑은 공의 0.90--0.92를 받지만 실제 몫은
0.18이고, 규칙 요리사는 $o$와 $s$가 비슷해 제대로 평가된다. 게임 단위로 보면 판정자의 공은 배달 비율과
$r=0.94$로 맞고(두 판정자 모두), Shapley 몫과는 $r=0.27$--$0.33$으로만 맞는다. 길막기의 음수 $s$를 0으로
잘라도 $r=0.54$--$0.56$에 그친다. 즉 판정자는 기여가 아니라 \emph{누가 마지막 단계를 했는지}로 공을 나눈다.
\TODO{GRPO 도둑, 그림자, 포기에 대한 판정 결과}

\textbf{반사실을 보지 못한다.} 도둑 옆에서 일한 두 준비 담당을 비교한다. 학습된 준비 담당은 혼자서는
아무것도 만들지 못하고($s=0.48$), 규칙 요리사는 혼자서도 모든 일을 할 수 있다($s=0.96$). 판정자는 둘에게
비슷한 공을 준다(GPT-4o 0.33 대 0.30, Qwen 0.47 대 0.43). 기록만으로는 누가 없어도 되는 팀원인지 알 수 없다.
프롬프트에 Shapley 정의를 주거나 혼자 했을 때의 결과까지 주어도, 예비 실험(6--7쌍)에서 Qwen은 도둑에게
여전히 0.68--0.70의 공을 준다. \TODO{정의/혼자 결과 프롬프트를 충분한 쌍으로}

% ===================================================================================== Limitations
\section{Limitations}
주요 결과는 \texttt{cramped\_room} 하나에서 나왔다. 두 팀원 모두 혼자서는 수프를 만들지 못하면 Shapley
몫이 공식상 0.5로 고정되어 정보가 없다. 멈춘 팀원도 부엌의 한 칸을 차지하므로 완전히 ``없는'' 상태는
아니다. seed는 3개이고 A는 seed 0 하나만 썼으며, 그림자와 포기는 greedy가 아니라 A 옆에서만 학습했다.
근접 보너스의 크기는 하나만 시험했다. LLM 판정자 실험의 pool은 greedy 옆의 다섯 종류뿐이며, 반사실 정보를
준 프롬프트는 아직 예비 결과이다.

% ===================================================================================== Conclusion
\section{Conclusion}
흔한 보상 설계는 팀 점수를 해치지 않거나 오히려 올리면서도 공을 부풀리는 나쁜 팀원을 만든다. 팀 점수
감독자는 점수를 낮추지 않는 이런 팀원을 통과시키고, LLM 판정자는 결과를 마무리한 팀원에게 공을 주어 이를
놓친다. 에이전트 팀에서 기여는 결과나 기록이 아니라 반사실로 측정해야 하며, 그래야 공을 공정하게 나눌 뿐
아니라 팀에 꼭 필요하지 않은 에이전트에 드는 비용도 찾아낼 수 있다. 우리의 나쁜 팀원 pool은 이런 credit
판정자를 시험하는 도구가 된다.

\bibliography{references}

% ===================================================================================== Appendix
\appendix
\section{Hyperparameters}
표~\ref{tab:hparams}는 모든 학습에 쓴 설정이다. 값은 \texttt{train/config.py}의 기본값이며, 각 결과 폴더의
\texttt{config.json}에 그대로 기록된다. 보상과 동료만 바뀌고 나머지는 모든 학습에서 같다.

\textbf{공통.} 정책은 $520 \rightarrow 128 \rightarrow 128 \rightarrow 6$ MLP(tanh)이고, 6개 행동(이동 4,
가만히, 상호작용)에 대한 categorical 분포를 낸다. 관측은 running mean/std로 정규화한다. 최적화는 Adam이다.
에피소드는 400 step에서 끊기므로, PPO는 마지막 보상을 $\gamma V(s_T)$로 bootstrap한다.

\textbf{PPO.} 4000 step(에피소드 10개)을 모은 뒤 GAE로 advantage를 계산하고, minibatch마다 정규화한다. 가치 함수는
정책과 별도의 MLP이다.

\textbf{GRPO.} 한 번의 업데이트마다 같은 reset seed에서 에피소드 8개를 현재 정책으로 돌린다. 그룹 전체의 step
보상을 평균과 표준편차로 정규화한 뒤, 각 step의 advantage를 그 뒤 정규화 보상의 할인 합으로 둔다
(\emph{process supervision}, \citealt{shao2024deepseekmath}). 이후 advantage를 표준편차 1로 맞추고, clip된
surrogate를 KL 항 없이 최적화한다. reset seed는 업데이트마다 바뀐다.

\begin{table}[h]
\centering
\small
\setlength{\tabcolsep}{4pt}
\begin{tabular}{@{}lcc@{}}
\toprule
 & PPO & GRPO \\
\midrule
\multicolumn{3}{@{}l}{\emph{환경과 보상}} \\
맵 & \multicolumn{2}{c}{\texttt{cramped\_room}} \\
에피소드 길이 & \multicolumn{2}{c}{400 step} \\
수프 하나 & \multicolumn{2}{c}{+20 (team 보상)} \\
줍기 보너스 (pickup) & \multicolumn{2}{c}{줍기마다 +3, 줄이지 않음} \\
근접 보너스 (proximity) & \multicolumn{2}{c}{동료 옆 step마다 +1, 줄이지 않음} \\
공식 shaping (team, individual) & \multicolumn{2}{c}{1M step까지 1 $\rightarrow$ 0 선형} \\
\midrule
\multicolumn{3}{@{}l}{\emph{학습}} \\
환경 step / 학습 & \multicolumn{2}{c}{$2\times10^6$} \\
seed & \multicolumn{2}{c}{0, 1, 2} \\
은닉층 & \multicolumn{2}{c}{2 $\times$ 128, tanh} \\
최적화 & \multicolumn{2}{c}{Adam} \\
학습률 & \multicolumn{2}{c}{$3\times10^{-4}$} \\
할인율 $\gamma$ & \multicolumn{2}{c}{0.99} \\
clip $\epsilon$ & \multicolumn{2}{c}{0.2} \\
gradient norm clip & \multicolumn{2}{c}{0.5} \\
epoch / 업데이트 & \multicolumn{2}{c}{8} \\
minibatch & \multicolumn{2}{c}{500} \\
업데이트당 데이터 & 4000 step & 8 에피소드 (3200 step) \\
advantage & GAE, $\lambda=0.98$ & 그룹 정규화 (process) \\
가치 함수 & MLP, 계수 0.5 & 없음 \\
entropy 계수 & 0.01 & 0 \\
KL 항 & -- & 없음 \\
\midrule
\multicolumn{3}{@{}l}{\emph{평가}} \\
평가 게임 & \multicolumn{2}{c}{10개, seed $10^4+k$} \\
행동 선택 & \multicolumn{2}{c}{정책에서 샘플링} \\
\bottomrule
\end{tabular}
\caption{하이퍼파라미터. PPO와 GRPO는 학습률, epoch, minibatch, clip을 같게 맞췄다.}
\label{tab:hparams}
\end{table}

\end{document}
